% !TeX root = ../thuthesis-example.tex

\clearpage
\thispagestyle{abstractcn}
\addcontentsline{toc}{chapter}{摘 要}

\vspace*{1.15cm}

\begin{center}
  {\heiti\bfseries\zihao{3} 摘\quad 要\par}
\end{center}

\vspace{0.25cm}

{\songti\zihao{-4}\setlength{\parindent}{2em}\linespread{1.65}\selectfont
本项目围绕生鲜水果零售与配送场景，设计并实现了一个基于 Spring Boot 与 Spring Cloud 的生鲜分布式运营系统。系统采用前后端分离与轻量微服务架构，管理端基于 RuoYi Vue3 与 Element Plus 实现，后端以 Spring Boot 为基础，结合 Nacos、Gateway、OpenFeign、Redis、RabbitMQ、MinIO、WebSocket 等技术完成服务治理、统一网关、服务调用、缓存、异步通知、对象存储和实时提醒等功能。

系统主要包括后台管理端、微信小程序端、统一网关、业务服务、支付服务和通知服务等模块。后台管理端支持水果分类管理、水果品种管理、订单管理和运营数据展示；微信小程序端支持用户浏览、下单与模拟支付；后端通过支付成功事件驱动订单状态更新和包装提醒推送，形成完整的业务闭环。系统本地运行验证采用 PM2 管理长运行进程，降低本地运行风险；同时保留容器化部署分支作为后续一键部署和工程化扩展方案。

通过本项目，能够展示 Spring Boot 单体业务开发能力、Spring Cloud 微服务治理能力、前后端分离开发能力以及分布式应用常见中间件的综合使用能力，满足分布式应用技术课程设计的实践要求。

\vspace{0.65cm}
\noindent{\heiti 关键词：}Spring Boot；Spring Cloud；Nacos；Gateway；OpenFeign；Redis；RabbitMQ；MinIO；WebSocket；Vue3
\par}

\clearpage
\thispagestyle{abstracten}
\addcontentsline{toc}{chapter}{Abstract}

\vspace*{0.9cm}

\begin{center}
  {\rmfamily\bfseries\zihao{3} Abstract\par}
\end{center}

\vspace{0.15cm}

{\rmfamily\zihao{-4}\setlength{\parindent}{2em}\linespread{1.45}\selectfont
This project designs and implements a distributed fresh-fruit retail and delivery operation system based on Spring Boot and Spring Cloud. The system adopts a front-end/back-end separated architecture and a lightweight microservice design. The management frontend is built with RuoYi Vue3 and Element Plus, while the backend integrates Nacos, Gateway, OpenFeign, Redis, RabbitMQ, MinIO, and WebSocket to support service governance, unified routing, inter-service calls, caching, asynchronous notification, object storage, and real-time order reminders.

The system includes an admin management frontend, a WeChat mini-program client, a gateway service, a business service, a payment service, and a notification service. The admin frontend supports fruit category management, fruit item management, order management, and operation dashboard display. The mini-program supports browsing, ordering, and mock payment. After payment succeeds, the backend updates the order status and pushes packaging reminders, forming a complete business workflow. For local validation, PM2 is used to manage long-running processes, while a deployment automation branch is reserved for one-click startup and engineering extension.

\vspace{0.45cm}
\noindent{\bfseries Keywords:} Spring Boot; Spring Cloud; Nacos; Gateway; OpenFeign; Redis; RabbitMQ; MinIO; WebSocket; Vue3
\par}
